\section{Introduction}\label{sec:introduction}

Turing reducibility compares the information available from sets of
integers: $A\leT B$ means that an oracle Turing machine with oracle $B$
computes the characteristic function of $A$. Mutual reducibility
defines Turing equivalence, and its equivalence classes form the
partial order $(\DT,\leT)$ of Turing degrees. The underlying notion
of relative computation goes back to Turing's oracle machines
\cite[Section~4]{Turing1939}; see \cite{Soare1987} for standard
degree-theoretic background. The automorphism problem asks how much
of this information the abstract ordering retains.

The rigidity conjecture for the Turing degrees asks whether every
automorphism of this ordering fixes every degree.
An automorphism here is any bijection
$\pi:\DT\to\DT$ satisfying
\[
\mathbf a\leT\mathbf b
\quad\Longleftrightarrow\quad
\pi(\mathbf a)\leT\pi(\mathbf b).
\]
We work in ZFC and impose no definability or regularity hypothesis
on $\pi$.

\begin{theorem}\label{thm:rigidity}
Every order automorphism of the full Turing-degree structure is
the identity:
\[
\forall\pi\in\Aut(\DT,\leT)\ \forall\mathbf a\in\DT
\qquad \pi(\mathbf a)=\mathbf a.
\]
\end{theorem}

\paragraph{Prior work and the remaining problem.}
The Turing jump sends a degree $\mathbf a$ to the degree
$\mathbf a'$ of the halting problem relative to an oracle of
degree $\mathbf a$. Jockusch and Solovay proved that every
jump-preserving order automorphism fixes all degrees above
$\mathbf0^{(4)}$, the fourth jump of the computable degree
\cite{JockuschSolovay1977}. Using Slaman and Woodin's definability
of the double jump, Shore and Slaman proved that the jump is
definable from the ordering
\cite{ShoreSlaman1999}, so every order automorphism preserves it.
Shore later gave direct degree-theoretic definitions of the jump
\cite{Shore2007}.

A cone consists of all degrees above a fixed degree. Nerode and
Shore proved that every automorphism fixes some cone, whose base
may depend on the automorphism \cite{NerodeShore1980}; see also
\cite[Theorem~3.2, author's version]{SlamanGlobal}.
Slaman and Woodin obtained a fixed cone common to all
automorphisms: every automorphism fixes all degrees above
$\mathbf0''$, the second jump of the computable degree.
Their analysis also shows that the automorphism group is
countable and that each automorphism has an arithmetic
representation on sets of integers
\cite{SlamanWoodin2005,SlamanGlobal}.
These results show that the order determines substantial
computability-theoretic structure. They leave the rigidity
problem at degrees outside the known fixed cone.

The representation theorem is the external input to our proof.
In its all-input form, it assigns to every arbitrary degree
automorphism a single arithmetic, hence Borel, function on sets of
integers that induces that automorphism on degrees
\cite[Theorem~6.3.1(2)]{SlamanWoodin2005}. The cited source is the
2005 unpublished manuscript; the result also appears in Slaman's
account \cite{SlamanGlobal}. This theorem supplies the regularity
needed for a category argument while retaining the unrestricted
scope of Theorem~\ref{thm:rigidity}.

\paragraph{A consequence for definability.}
Theorem~\ref{thm:rigidity} also settles the biinterpretability
conjecture of Slaman and Woodin. Full second-order arithmetic is
the two-sorted structure of natural numbers and all subsets of
natural numbers, with arithmetic and membership. Slaman and Woodin
interpret this structure using finite tuples of Turing degrees.
In their formulation, biinterpretability asks whether the relation
matching such a code for a set $X\subseteq\N$ with its degree
$\degT X$ is first-order definable in the degree ordering without
parameters \cite[Definition~7.5.1]{SlamanWoodin2005}.
They prove that this is equivalent to rigidity
\cite[Theorem~7.5.3]{SlamanWoodin2005}; thus our theorem gives a
positive answer. We state the resulting corollary in
Section~\ref{sec:rigidity}.

\paragraph{Recovery and removal of parameters.}
Our main construction is the recovery theorem of
Proposition~\ref{prop:recovery}. We identify $2^\N$ with the sets of
integers; recovering a real means computing which rational numbers
lie below it. Write $\oplus$ for the oracle join,
which combines the information in finitely many sets of integers.
For a Borel map $F:\R\to2^\N$ with countable fibers satisfying
\[
F(x+y)\leT F(x)\oplus F(y),
\]
we recover any prescribed irrational $t\in(0,1)$ from four values
of $F$ at rational affine combinations of $t$ and two auxiliary
reals $u,v$. Recovery holds for a comeager set of pairs $(u,v)$:
the exceptional pairs form a countable union of nowhere-dense
sets. The statement applies beyond functions representing degree
automorphisms.

The numerical mechanism is an averaging identity. For a fixed
positive rational $\delta$, start with a real state $s_0=0$
and repeatedly choose the increment $\delta t$ or
$\delta(t-1)$, recording the choice as $r_n=0$ or $r_n=1$,
respectively. Summing these increments gives
\[
s_N=\delta\left(Nt-\sum_{n<N}r_n\right).
\]
If the states remain in $(-1,1)$, the proportion of choices with
$r_n=1$ approximates $t$ with error less than $1/(N\delta)$.
Thus computing the binary choices suffices to recover the
irrational parameter $t$ with an effective error bound.

The function $F$ supplies both a rule for choosing the increments
and a way to compute those choices. Countable fibers ensure that
$F$ takes distinct values on opposite sides of $u$. Continuity
on a comeager restriction then makes one output bit take opposite
values near two such points. At state $s_n$, we test this bit of
$F(u+s_n)$, choosing an upward step near the lower endpoint and
a downward step near the upper endpoint. This keeps the states
bounded, regardless of the bit's behavior between the two
endpoint regions.

The addition reductions need not be uniform in their inputs.
For each auxiliary pair in a comeager set, Borel regularity
provides two programs that remain fixed along the recurrence.
The first adds the selected increment while moving from a
neighborhood of $u$ to a neighborhood of $v$; the second returns
to the neighborhood of $u$ at the new offset. The computation uses a
fixed function value for each of the two possible increments,
a fixed return value, and the initial value $F(u)$.
These four values are its only real oracles. The program indices
and rational constants may depend on $t$ and the chosen pair
$(u,v)$.

Finite-oracle recovery also appears in perfect-set coding.
Groszek and Slaman recover a real from two branches of a perfect
tree together with a sequence supplying a suitable dense set
\cite[Lemma~2.2, author's version]{GroszekSlaman1998}.
Lutz and Siskind's refinement
recovers any real from four elements of a perfect set together
with a fixed real that computes each member of a countable dense
subset, without requiring an enumeration computable from that
fixed real
\cite[Theorem~4.3, arXiv version]{LutzSiskind2025}.
Our recovery statement has different hypotheses: it uses the
addition relation for one Borel function and prescribes affine
arguments for its four values. Its bounded recurrence supplies
the decoding procedure without an additional real oracle.

To conclude rigidity, represent an arbitrary automorphism by a
map $F$ as above. The degrees of the four arguments are bounded by
the join of the degrees of $t,u,v$. Applying the inverse
automorphism to the recovery bound therefore bounds the preimage
degree of $t$ by this same join. Sacks's relative category
cone-avoidance principle says that, over a fixed oracle,
auxiliary reals compute a fixed set not already computable from
that oracle only on a meager set of choices
\cite[Lemma~4.1]{KumarShelah2023}. Since recovery holds on a
comeager set, Lemma~\ref{lem:avoidance} removes the auxiliary
degrees and yields $\pi^{-1}(\mathbf a)\leT\mathbf a$ for every
degree $\mathbf a$. Applying $\pi$ to this inequality and repeating
the argument for $\pi^{-1}$ gives the opposite inequalities between
$\pi(\mathbf a)$ and $\mathbf a$.

A related method appears in Kjos-Hanssen's proof of rigidity for
automorphisms induced by permutations of the natural numbers
\cite[Theorem~20, arXiv version]{KjosHanssen2018}.
That argument recovers a Bernoulli parameter (the probability of
a $1$), uses a measure cone theorem, and applies the inverse
automorphism. Here the
parameter is recovered from the addition relation for a Borel
function, and category removes the auxiliary degrees. The
representing function need not arise from a permutation.

\paragraph{Organization.}
Section~\ref{sec:preliminaries} fixes the real and oracle codings
and proves the category lemmas. Section~\ref{sec:representation}
states the exact Slaman--Woodin input and constructs the Borel map
on the real line. Section~\ref{sec:recovery} proves four-value
recovery, and Section~\ref{sec:rigidity} removes the auxiliary
degrees, completes the proof, and derives the biinterpretability
corollary.
